%%%%%%%%%%%%%%%%%%%%%%%%%%%%%%%%%%%%%%%%%%%%%%%%%%%%%%%%%%%%%%%%%%%%%%%%%%%
% Chapter: Voltage References
% Falstad examples: examples/references (generated by falstad/circuits/references.py)
%%%%%%%%%%%%%%%%%%%%%%%%%%%%%%%%%%%%%%%%%%%%%%%%%%%%%%%%%%%%%%%%%%%%%%%%%%%

%%%%%%%%%%%%%%%%%%%%%%%%%%%%%%%%%%%%%%%%%%%%%%%%%%%%%%%%%%%%%%%%%%%%%%%%%%%
\section{Why Voltage References?}
%%%%%%%%%%%%%%%%%%%%%%%%%%%%%%%%%%%%%%%%%%%%%%%%%%%%%%%%%%%%%%%%%%%%%%%%%%%
\nsl{Every measurement compares an unknown quantity with a known one. In an
electronic measurement chain the known quantity is almost always a
voltage: the \emph{reference voltage}. An analog-to-digital converter does
not measure volts, it measures the ratio of its input voltage to its
reference voltage. A bridge circuit does not deliver millivolts per kelvin,
it delivers a fraction of its excitation voltage. If the reference is
\SI{1}{\percent} too high, the result is \SI{1}{\percent} wrong, no matter
how good the rest of the circuit is.
\newline
A supply voltage from a voltage regulator is not good enough for this job.
A typical regulator has an initial tolerance of \SIrange{2}{5}{\percent},
changes with load and temperature and carries ripple and noise from the
digital part of the circuit. A voltage reference is a circuit (usually one
small chip) that is optimised for a single purpose: a voltage that is
accurate, stable over temperature and time, and quiet. It is usually not
meant to deliver much current.
\newline
In this chapter we look at the specifications of references, at the two
physical principles behind almost all of them (the Zener diode and the
bandgap of silicon), at the shunt reference TL431 and at the practical
questions: how to buffer and scale a reference, how to drive the reference
ladder of a flash converter, and how a \emph{ratiometric} design removes
the reference error from the result altogether. The capstone task needs
exactly this: a \SI{4.00}{\volt} reference for the ADC ladder and for the
sensor bridge \cite{horowitz2015,kester2005}.\newpage}

\begin{description}
\item[{\rot\bf The reference sets the scale of the whole chain}]~\\[-\lblskip]
\begin{center}
\begin{tikzpicture}[font=\sffamily\small, >=Stealth,
	box/.style={draw,thick,minimum height=9mm,minimum width=18mm,align=center}]
	\node[box,fill=rwuviolet!12] (sen) at (0,0) {sensor\\bridge};
	\node[box,fill=rwucyan!10] (amp) at (3,0) {amplifier};
	\node[box,fill=rwucyan!10] (adc) at (6,0) {ADC};
	\node[box,fill=rwuviolet!25] (ref) at (3,-2) {voltage\\reference};
	\draw[->,thick] (sen) -- (amp);
	\draw[->,thick] (amp) -- (adc);
	\draw[->,thick] (adc) -- ++(1.6,0) node[right]{code};
	\draw[->,thick] (ref) -| node[pos=0.75,left]{excitation} (sen);
	\draw[->,thick] (ref) -| node[pos=0.75,right]{$\Vref$} (adc);
\end{tikzpicture}
\end{center}
\bi
	\ite ADC: $\mathrm{code} = 2^N\cdot \Vin/\Vref$ \quad (an ADC measures a \emph{ratio})
	\ite bridge: output is a fraction of the excitation voltage
	\ite an error of $\Vref$ is a \textbf{gain error} of the complete chain
\ei
\os{\newpage}

\item[{\rot\bf Regulator or reference?}]~\\[-\lblskip]
\begin{center}
\begin{tabular}{lcc}
\toprule
 & voltage regulator & voltage reference \\
\midrule
purpose & supply power & define a voltage \\
initial accuracy & \SIrange{1}{5}{\percent} & \SIrange{0.02}{1}{\percent} \\
temperature coefficient & \SIrange{100}{300}{ppm/K} & \SIrange{1}{50}{ppm/K} \\
output current & \si{\milli\ampere} \ldots\ \si{\ampere} & \si{\micro\ampere} \ldots\ \SI{10}{\milli\ampere} \\
noise & high & low, specified \\
\bottomrule
\end{tabular}
\end{center}
\bi
	\ite a reference is a ``voltage standard'' on the board, not a power supply
	\ite if a load needs current: buffer the reference with an op-amp
\ei
\os{\newpage}
\end{description}

%%%%%%%%%%%%%%%%%%%%%%%%%%%%%%%%%%%%%%%%%%%%%%%%%%%%%%%%%%%%%%%%%%%%%%%%%%%
\section{Specifications}
%%%%%%%%%%%%%%%%%%%%%%%%%%%%%%%%%%%%%%%%%%%%%%%%%%%%%%%%%%%%%%%%%%%%%%%%%%%
\nsl{Data sheets of voltage references list a handful of parameters. They
all describe how far the actual output voltage can be from its nominal
value, but for different reasons: production spread, temperature,
supply voltage, load current, noise and ageing. For an error budget we add
them up (see the end of this chapter). The temperature coefficient is
usually given in ppm/K (parts per million per kelvin):
\SI{1}{ppm} $= 10^{-6}$, so \SI{100}{ppm} $= \SI{0.01}{\percent}$.\newpage}

\begin{description}
\item[{\rot\bf Data sheet parameters}]~\\[-\lblskip]
\bi
	\ite \textbf{initial accuracy}: tolerance at \SI{25}{\celsius}, e.g.\ $\pm\SI{0.1}{\percent}$
	\ite \textbf{temperature coefficient} (TC): e.g.\ \SI{20}{ppm/K}
	\ite \textbf{line regulation}: $\Delta\Vout/\Delta V_\mathrm{in}$, e.g.\ \SI{50}{\micro\volt/\volt}
	\ite \textbf{load regulation}: $\Delta\Vout/\Delta I_\mathrm{out}$, e.g.\ \SI{100}{ppm/\milli\ampere}
	\ite \textbf{noise}: \SIrange{0.1}{10}{\hertz} peak-to-peak, e.g.\ \SI{10}{\micro\volt}\,pp
	\ite \textbf{long-term drift}: e.g.\ \SI{50}{ppm} per \SI{1000}{\hour}
	\ite \textbf{thermal hysteresis}: change after a temperature cycle
\ei
\nsl{The low-frequency noise from \SIrange{0.1}{10}{\hertz} is given
separately because it cannot be filtered: it looks like a slow drift of
the reference. Broadband noise above that can be reduced with a capacitor
at a noise-reduction pin or with an RC filter after the reference.}
\os{\newpage}

\item[{\rot\bf Temperature coefficient: the box method}]~\\[-\lblskip]
\begin{center}
\begin{tikzpicture}
\begin{axis}[width=0.62\textwidth,height=42mm,xmin=-40,xmax=125,ymin=2.4972,ymax=2.5012,
	xlabel={$T$ / \si{\celsius}},ylabel={$\Vout$ / V},axis lines=left,
	yticklabel style={/pgf/number format/fixed,/pgf/number format/precision=4},
	font=\scriptsize,ytick={2.498,2.499,2.500,2.501}]
	\addplot[rwuviolet,thick,samples=60,domain=-40:125] {2.5007 - 4.0e-7*(x-35)^2 + 2e-6*(x-35)};
	\draw[rwucyandark,dashed] (axis cs:-40,2.49764) rectangle (axis cs:125,2.50070);
\end{axis}
\end{tikzpicture}
\end{center}
\keyeq{Box method}{$\displaystyle TC = \frac{V_\mathrm{max}-V_\mathrm{min}}{V_\mathrm{nom}\,(T_\mathrm{max}-T_\mathrm{min})}\cdot 10^6\ \mathrm{ppm}$}
\nsl{The curve in the figure is typical for a bandgap reference: it is a
parabola with its maximum near room temperature. The box method draws the
smallest rectangle around the curve over the specified temperature range
and divides its height by its width. In the example the box is
\SI{3.06}{\milli\volt} high and \SI{165}{\kelvin} wide, so
$TC = \SI{3.06}{\milli\volt}/(\SI{2.5}{\volt}\cdot\SI{165}{\kelvin}) \approx
\SI{7.4}{ppm/K}$. The TC does \emph{not} mean that the voltage changes
linearly; it only guarantees the box.}
\os{\newpage}
\end{description}

%%%%%%%%%%%%%%%%%%%%%%%%%%%%%%%%%%%%%%%%%%%%%%%%%%%%%%%%%%%%%%%%%%%%%%%%%%%
\section{The Zener Reference}
%%%%%%%%%%%%%%%%%%%%%%%%%%%%%%%%%%%%%%%%%%%%%%%%%%%%%%%%%%%%%%%%%%%%%%%%%%%
\nsl{The oldest and simplest reference is a Zener diode operated in reverse
breakdown. Above the breakdown voltage the current rises steeply, so the
voltage stays nearly constant over a wide range of currents. Two physical
effects cause the breakdown. Below about \SI{5}{\volt} it is the
\emph{Zener effect} (quantum mechanical tunnelling through a thin, heavily
doped junction); its breakdown voltage falls with temperature. Above about
\SI{6}{\volt} it is the \emph{avalanche effect} (carriers gain enough
energy to ionise further atoms); its breakdown voltage rises with
temperature. Between \SI{5}{\volt} and \SI{6}{\volt} both effects are
present and their temperature coefficients nearly cancel. This is why
\SI{5.6}{\volt} and \SI{6.2}{\volt} Zener diodes are preferred as
references \cite{tietze2008}.\newpage}

\begin{description}
\item[{\rot\bf Temperature coefficient of Zener diodes}]~\\[-\lblskip]
\begin{center}
\begin{tikzpicture}
\begin{axis}[width=0.7\textwidth,height=44mm,xmin=2.5,xmax=12.5,ymin=-0.09,ymax=0.09,
	xlabel={Zener voltage $V_Z$ / V},ylabel={TC / (\%/K)},axis lines=left,
	font=\scriptsize,grid=major,
	yticklabel style={/pgf/number format/fixed,/pgf/number format/precision=2}]
	\addplot[rwuviolet,thick,smooth,mark=*,mark size=1.2pt] coordinates {
		(2.7,-0.075) (3.3,-0.065) (3.9,-0.050) (4.7,-0.020) (5.1,-0.005)
		(5.6,0.015) (6.2,0.035) (6.8,0.045) (8.2,0.060) (10,0.070) (12,0.077)};
	\node[font=\scriptsize,align=center] at (axis cs:3.8,0.05) {Zener effect\\TC $<0$};
	\node[font=\scriptsize,align=center] at (axis cs:10,-0.05) {avalanche effect\\TC $>0$};
\end{axis}
\end{tikzpicture}
\end{center}
\bi
	\ite minimum $|TC|$ near \SIrange{5}{6}{\volt}: preferred reference Zeners \SI{5.6}{\volt}, \SI{6.2}{\volt}
	\ite \SI{6.2}{\volt} Zener ($+\SI{2}{\milli\volt/\kelvin}$) in series with a forward diode ($-\SI{2}{\milli\volt/\kelvin}$): compensated reference diode (1N821 family)
\ei
\nsl{The curve shows typical values of small-signal Zener diodes; the
spread between data sheets is large. The compensated reference diodes
achieve \SIrange{5}{100}{ppm/K}, but only at the specified current
(typically \SI{7.5}{\milli\ampere}). The best references ever built
(LM399, LTZ1000) use a \emph{buried} Zener below the chip surface, which
is far less noisy, together with an on-chip heater that keeps the die at a
constant temperature.}
\os{\newpage}

\item[{\rot\bf Feeding a Zener through a resistor}]~\\[-\lblskip]
\begin{center}
\begin{circuitikz}[scale=1.0, transform shape]
  \draw (0,0) to[V, invert, l=$V_S$] (0,3) to[R, l=$R_V$, i=$I$] (3,3)
        to[zD, invert, l_=$V_Z$] (3,0) -- (0,0);
  \draw (3,3) -- (5,3) to[R, l=$R_L$, i>^=$I_L$] (5,0) -- (3,0);
  \draw (3,3) to[short,*-] (3,3);
  \draw (0,0) node[ground]{};
  \draw (6,3) node[right]{$I_Z = I - I_L$};
\end{circuitikz}
\end{center}
\bi
	\ite $\displaystyle R_V = \frac{V_S - V_Z}{I_Z + I_L}$, \quad $I_Z$ must stay above $I_{Z,\min}$ for all $V_S$ and $I_L$
	\ite dynamic resistance $r_Z = \mathrm{d}V_Z/\mathrm{d}I_Z$: \SIrange{5}{50}{\ohm}
	\ite line regulation $\displaystyle\frac{\Delta V_Z}{\Delta V_S} \approx \frac{r_Z}{R_V + r_Z}$, \quad load regulation $\Delta V_Z \approx -r_Z\,\Delta I_L$
\ei
\falstad{zener_resistor}
\nsl{\par The Zener and the series resistor form a voltage divider whose lower
part has the small dynamic resistance $r_Z$. A supply change is divided
down by $r_Z/(R_V+r_Z)$, typically 1:100. The load current is taken from
the Zener current: if the load draws more, the Zener carries less, and its
voltage drops by $r_Z\,\Delta I_L$. Falstad models the Zener with an ideal
exponential characteristic $I_Z = \SI{5}{\milli\ampere}\cdot
e^{(V-V_Z)/V_T}$, so the nominal voltage is reached at
\SI{5}{\milli\ampere} and the dynamic resistance is only $V_T/I_Z \approx
\SI{5}{\ohm}$. Real \SI{5.6}{\volt} Zeners have \SIrange{10}{40}{\ohm}, so
in reality the regulation is several times worse than in the simulation.}
\os{\newpage}

\item[{\rot\bf Feeding a Zener from a current source}]~\\[-\lblskip]
\begin{center}
\begin{circuitikz}[scale=1.0, transform shape]
  \draw (0,0) to[V, invert, l=$V_S$] (0,3) -- (2.5,3)
        to[I, l=$I_0$] (2.5,1.5) to[short,-*] (2.5,1.5)
        to[zD, invert, l_=$V_Z$] (2.5,0) -- (0,0);
  \draw (2.5,3) -- (4,3) to[R, l=$r_o$] (4,1.5) -- (2.5,1.5);
  \draw (2.5,1.5) to[short,-o] (5.5,1.5) node[right]{$\Vref$};
  \draw (0,0) node[ground]{};
\end{circuitikz}
\end{center}
\bi
	\ite the Zener current no longer depends on $V_S$
	\ite only the output resistance $r_o$ of a real current source remains:
	$\displaystyle\frac{\Delta V_Z}{\Delta V_S} \approx \frac{r_Z}{r_o}$
	\ite e.g.\ transistor current source from chapter~2, $r_o = \SI{100}{\kilo\ohm}$
\ei
\falstad{zener_current_source}
\nsl{\par With $r_Z = \SI{5}{\ohm}$ and $r_o = \SI{100}{\kilo\ohm}$ the line
regulation improves from $1:240$ (with $R_V = \SI{1.2}{\kilo\ohm}$) to
$1:20000$. The same trick is used inside every integrated reference: the
reference element is fed from a current that is derived from the reference
voltage itself, so that supply changes do not reach it.}
\os{\newpage}
\end{description}

\exercise{Zener reference with series resistor}
A Zener diode with $V_Z = \SI{5.6}{\volt}$ (at \SI{5}{\milli\ampere}) is fed
from $V_S = \SI{12}{\volt}$ through $R_V = \SI{1.2}{\kilo\ohm}$. Use the
Falstad model of the Zener, $I_Z = \SI{5}{\milli\ampere}\cdot e^{(V-V_Z)/V_T}$
with $V_T = \SI{25.9}{\milli\volt}$.
\begin{talist}
\ta Calculate the Zener current and the dynamic resistance $r_Z$.
\ta How much does $\Vref$ change if $V_S$ rises to \SI{14}{\volt}?
\ta A load $R_L = \SI{2.2}{\kilo\ohm}$ is connected at \SI{12}{\volt}. How much does $\Vref$ drop?
\end{talist}

\begin{solution}
\begin{talist}
\ta $I_Z \approx (\SI{12}{\volt}-\SI{5.6}{\volt})/\SI{1.2}{\kilo\ohm} = \SI{5.33}{\milli\ampere}$,
	so $\Vref = V_Z + V_T\ln(5.33/5) = \SI{5.602}{\volt}$ and
	$r_Z = V_T/I_Z = \SI{25.9}{\milli\volt}/\SI{5.33}{\milli\ampere} = \SI{4.9}{\ohm}$.
\ta At \SI{14}{\volt}: $I_Z = \SI{7.0}{\milli\ampere}$,
	$\Delta\Vref = V_T \ln(7.0/5.33) = \SI{7.0}{\milli\volt}$
	(small-signal estimate with the mean $r_Z \approx \SI{4.2}{\ohm}$:
	$\SI{2}{\volt}\cdot 4.2/1204 = \SI{7.0}{\milli\volt}$).
	Line regulation \SI{3.5}{\milli\volt/\volt}, i.e.\ \SI{0.06}{\percent/\volt}.
\ta The load takes $I_L = \SI{5.6}{\volt}/\SI{2.2}{\kilo\ohm} = \SI{2.55}{\milli\ampere}$,
	the Zener keeps $I_Z = \SI{5.33}{}-\SI{2.55}{\milli\ampere} = \SI{2.79}{\milli\ampere}$:
	$\Delta\Vref = V_T\ln(2.79/5.33) = -\SI{16.6}{\milli\volt}$.
\item[] \falstad{zener_resistor}: the labels show \texttt{Vz12} $=\SI{5.602}{\volt}$,
	\texttt{Vz14} $=\SI{5.609}{\volt}$ ($+\SI{7.0}{\milli\volt}$) and
	\texttt{Vz12L} $=\SI{5.585}{\volt}$ ($-\SI{16.6}{\milli\volt}$).
\end{talist}
\end{solution}

\exercise{Zener with current-source feed}
The Zener of the previous exercise is fed from a current source of
\SI{5}{\milli\ampere} with an output resistance $r_o=\SI{100}{\kilo\ohm}$.
How much does $\Vref$ change when $V_S$ rises from \SI{12}{\volt} to
\SI{14}{\volt}? Compare with the series resistor.

\begin{solution}
\begin{talist}
\ta The current through $r_o$ changes by $\SI{2}{\volt}/\SI{100}{\kilo\ohm} = \SI{20}{\micro\ampere}$,
	the Zener current from \SI{5.064}{\milli\ampere} to \SI{5.084}{\milli\ampere}.
	$\Delta\Vref = V_T\ln(5.084/5.064) = \SI{0.10}{\milli\volt}$,
	70 times less than with $R_V = \SI{1.2}{\kilo\ohm}$ (\SI{7.0}{\milli\volt}).
\item[] \falstad{zener_current_source}: at \SI{12}{\volt} \texttt{VzI} $=\SI{5.600}{\volt}$;
	moving the slider to \SI{14}{\volt} changes \texttt{VzR} by \SI{7.0}{\milli\volt}
	and \texttt{VzI} by only \SI{0.1}{\milli\volt}.
\end{talist}
\end{solution}

%%%%%%%%%%%%%%%%%%%%%%%%%%%%%%%%%%%%%%%%%%%%%%%%%%%%%%%%%%%%%%%%%%%%%%%%%%%
\section{The Bandgap Reference}
%%%%%%%%%%%%%%%%%%%%%%%%%%%%%%%%%%%%%%%%%%%%%%%%%%%%%%%%%%%%%%%%%%%%%%%%%%%
\nsl{Zener diodes need more than \SI{6}{\volt} of supply and several
milliamperes. Modern circuits run from \SI{3.3}{\volt} or less and from
microamperes. Almost all integrated references (and the references inside
microcontroller ADCs) therefore use the \emph{bandgap} principle, invented
by Widlar \cite{widlar1971} and brought to its classic form by Brokaw
\cite{brokaw1974}. The idea: the base-emitter voltage of a transistor falls
with temperature, and the difference of two base-emitter voltages at
different current densities rises with temperature. A suitably weighted
sum of the two does not depend on temperature. The sum turns out to be
close to the bandgap voltage of silicon, \SI{1.2}{\volt}, which gives the
circuit its name.\newpage}

\begin{description}
\item[{\rot\bf Two voltages with opposite temperature coefficients}]~\\[-\lblskip]
\bi
	\ite base-emitter voltage: $V_{BE} = V_T\ln(I_C/I_S)$ with $V_T = kT/q = \SI{25.7}{\milli\volt}$ at \SI{25}{\celsius}
	\itee $I_S$ rises strongly with $T$: $V_{BE}$ falls by about \SI{2}{\milli\volt/\kelvin} (CTAT: complementary to absolute temperature)
	\ite two transistors with current densities in the ratio $N$:
	\keyeq{Delta-VBE}{$\displaystyle \Delta V_{BE} = V_{BE1}-V_{BE2} = V_T\ln N = \frac{kT}{q}\ln N$ \qquad (PTAT: proportional to absolute temperature)}
	\itee $k/q = \SI{86.2}{\micro\volt/\kelvin}$; $N=8$: $\Delta V_{BE} = \SI{53.5}{\milli\volt}$ at \SI{25}{\celsius}, $+\SI{0.179}{\milli\volt/\kelvin}$
	\itee $I_S$ cancels: $\Delta V_{BE}$ depends only on $T$ and $N$
\ei
\os{\newpage}

\item[{\rot\bf Delta-VBE in Falstad}]~\\[-\lblskip]
\begin{center}
\begin{circuitikz}[scale=0.9, transform shape]
  \draw (-0.6,4) -- (4.6,4);
  \draw (0,4) to[I, l_=\SI{100}{\micro\ampere}] (0,2.4) -- (0,2.2) node[npn,anchor=C](Q1){};
  \draw (Q1.B) -- ++(-0.3,0) |- (0,2.3);
  \draw (Q1.E) node[ground]{};
  \draw (0,2.3) node[circ]{} -- (0.6,2.3) node[right]{$V_{BE1}$};
  \draw (4,4) to[I, l_=\SI{100}{\micro\ampere}] (4,2.4) -- (4,2.2) node[npn,anchor=C](Q2){};
  \draw (Q2.B) -- ++(-0.3,0) |- (4,2.3);
  \draw (Q2.E) node[ground]{};
  \draw (4,2.3) node[circ]{} -- (4.6,2.3) node[right]{$V_{BE8}$};
  \node[right,font=\small] at (4.3,1.2) {$8\times$ area};
  \node[right,font=\small] at (0.3,1.2) {$1\times$ area};
\end{circuitikz}
\end{center}
\bi
	\ite same current, 8 transistors in parallel: current density $1/8$
	\ite Falstad has no temperature, but $\Delta V_{BE}$ at room temperature can be measured:
	$V_{BE1} = \SI{0.536}{\volt}$, $V_{BE1}-V_{BE8} = \SI{53.8}{\milli\volt}$
	\ite current ratio 10: $\Delta V_{BE} = V_T \ln 10 = \SI{59.6}{\milli\volt}$ (``\SI{60}{\milli\volt} per decade'')
\ei
\falstad{delta_vbe}
\nsl{\par Falstad uses $V_T = \SI{25.865}{\milli\volt}$ (\SI{27}{\celsius}) and a
transistor model with $I_S = \SI{e-13}{\ampere}$. The absolute value of
$V_{BE}$ is therefore lower than in a real small-signal transistor
(about \SI{0.65}{\volt} at \SI{100}{\micro\ampere}), but the difference
$V_T\ln N$ is exact, because $I_S$ cancels. In an integrated circuit the
area ratio is made by placing $N$ identical transistors in parallel, exactly
as in the simulation; this is the most precise ratio a chip can offer.}
\os{\newpage}

\item[{\rot\bf The bandgap principle}]~\\[-\lblskip]
\begin{center}
\begin{tikzpicture}
\begin{axis}[width=0.72\textwidth,height=46mm,xmin=-40,xmax=125,ymin=0,ymax=1.45,
	xlabel={$T$ / \si{\celsius}},ylabel={V},axis lines=left,font=\scriptsize,
	legend style={at={(1.02,0.5)},anchor=west,draw=none,font=\scriptsize},
	declare function={TK(\t)=\t+273.15; VBE(\t)=1.205-(1.205-0.65)*TK(\t)/298.15-2*8.617e-5*TK(\t)*ln(TK(\t)/298.15); PT(\t)=0.6063*TK(\t)/298.15;}]
	\addplot[rwucyan,thick,domain=-40:125,samples=40] {VBE(x)};
	\addplot[rwucyandark,thick,dashed,domain=-40:125,samples=40] {PT(x)};
	\addplot[rwuviolet,very thick,domain=-40:125,samples=40] {VBE(x)+PT(x)};
	\legend{$V_{BE}$ ($-\SI{2}{\milli\volt/\kelvin}$),$K\,\Delta V_{BE}$ (PTAT),sum $\approx\SI{1.25}{\volt}$}
\end{axis}
\end{tikzpicture}
\end{center}
\keyeq{Bandgap reference}{$\displaystyle V_{BG} = V_{BE} + K\cdot V_T\ln N \approx \SI{1.25}{\volt}$, \qquad $K$ chosen so that $\mathrm{d}V_{BG}/\mathrm{d}T = 0$ at room temperature}
\nsl{The PTAT voltage must rise by \SI{2}{\milli\volt/\kelvin}. With $N=8$
one $\Delta V_{BE}$ rises by \SI{0.179}{\milli\volt/\kelvin}, so it must be
amplified by $K\approx 11$. At room temperature $K\,\Delta V_{BE} \approx
11\cdot\SI{54}{\milli\volt} \approx \SI{0.6}{\volt}$, and the sum is
\SIrange{1.2}{1.25}{\volt}. Extrapolated to \SI{0}{\kelvin}, the sum is
exactly the bandgap voltage of silicon, $V_{G0} = \SI{1.205}{\volt}$. The
remaining curvature (the $T\ln T$ term in $V_{BE}$) gives the parabola
seen in the box-method figure; simple bandgaps reach
\SIrange{10}{50}{ppm/K}, curvature-corrected ones \SIrange{1}{5}{ppm/K}.}
\os{\newpage}

\item[{\rot\bf The Brokaw cell}]~\\[-\lblskip]
\begin{center}
\begin{circuitikz}[scale=0.95, transform shape]
  \draw (0.4,7.4) node[left]{$+V_S$} -- (4,7.4);
  \draw (1,7.4) to[R, l_=$R_C$] (1,5.8) -- (1,4.0) node[npn,anchor=C,xscale=-1](Q2){};
  \draw (4,7.4) to[R, l=$R_C$] (4,5.8) -- (4,4.0) node[npn,anchor=C](Q1){};
  \node[font=\small] at (-0.1,3.2) {Q2: $1\times$};
  \node[font=\small] at (5.1,3.2) {Q1: $8\times$};
  \draw (Q2.B) -- (Q1.B);
  \draw (Q2.E) -- (1,1.0) -- (2.5,1.0);
  \draw (Q1.E) to[R, l=$R_2$] (4,1.0) -- (2.5,1.0) node[circ]{} node[above left]{$V_E$}
        to[R, l=$R_1$] (2.5,-0.6) node[ground]{};
  \node[op amp, noinv input up] (A) at (7.0,4.9) {};
  \draw (1,4.8) node[circ]{} -- (2,4.8) |- (A.+);
  \draw (4,4.4) node[circ]{} -- (5.3,4.4) |- (A.-);
  \draw (Q2.B -| 2.5,0) node[circ]{} -- (2.5,2.4) -- (8.6,2.4);
  \draw (A.out) -- (8.6,4.9) -- (8.6,2.4) node[circ]{} to[short,-o] (9.4,2.4) node[right]{$V_{BG}$};
\end{circuitikz}
\end{center}
\nsl{Q1 consists of eight transistors in parallel, Q2 of one. Both bases
are driven by the output $V_{BG}$. The op-amp compares the voltages at the
two equal collector resistors $R_C$ and adjusts $V_{BG}$ until both
collector currents are equal. Then the current density in Q2 is eight times
that in Q1, and the difference of the base-emitter voltages appears across
$R_2$.}
\os{\newpage}

\item[{\rot\bf Brokaw cell: how it works}]~\\[-\lblskip]
\bi
	\ite the op-amp forces equal collector currents $I$ ($R_C$ equal, inputs at the collectors)
	\ite current density ratio $N = 8$: $\Delta V_{BE} = V_T\ln 8$ appears across $R_2$
	\itee $I = V_T\ln 8 / R_2$ \quad (PTAT current)
	\ite $R_1$ carries both emitter currents: $V_E = 2 I R_1 = 2\,\dfrac{R_1}{R_2}\,V_T\ln 8$
	\ite output: \quad $V_{BG} = V_{BE2} + 2\dfrac{R_1}{R_2}\,V_T \ln 8$ \quad ($K = 2R_1/R_2$)
	\ite Falstad, $R_2 = \SI{1}{\kilo\ohm}$, $R_1 = \SI{5.6}{\kilo\ohm}$:
	$I = \SI{53.8}{\micro\ampere}$, $V_E = \SI{0.602}{\volt}$, $V_{BE2} = \SI{0.520}{\volt}$, $V_{BG} = \SI{1.122}{\volt}$
\ei
\falstad{brokaw_cell}
\nsl{\par With $K = 2\cdot 5.6 = 11.2$ the PTAT part rises by
$11.2\cdot\SI{0.179}{\milli\volt/\kelvin} = \SI{2.0}{\milli\volt/\kelvin}$
and compensates $V_{BE}$. With a real transistor ($V_{BE}\approx
\SI{0.65}{\volt}$) the output would be \SI{1.25}{\volt}. Falstad's transistor
model has a lower $V_{BE}$ and no temperature dependence, so it shows
\SI{1.122}{\volt}; move the slider of $R_1$ to see how the PTAT part
scales ($R_1 = \SI{6.8}{\kilo\ohm}$ gives \SI{1.25}{\volt}). In a real
chip $R_1$ is trimmed at the wafer test until the output has zero TC. A
resistive divider between $V_{BG}$ and the op-amp output (not shown) gives
larger outputs such as \SI{2.5}{\volt} or \SI{4.096}{\volt}.}
\os{\newpage}
\end{description}

\exercise{Bandgap design}
A Brokaw cell has an area ratio $N = 8$, $R_2 = \SI{1}{\kilo\ohm}$ and
$R_1 = \SI{5.6}{\kilo\ohm}$. Use $V_T = \SI{25.865}{\milli\volt}$ (as Falstad).
\begin{talist}
\ta Calculate $\Delta V_{BE}$, the collector current $I$ and $V_E$.
\ta The Falstad transistor model gives $V_{BE2} = \SI{0.520}{\volt}$ at this current. Calculate $V_{BG}$.
\ta Calculate the temperature coefficient of the PTAT part and compare it with $\mathrm{d}V_{BE}/\mathrm{d}T \approx -\SI{2}{\milli\volt/\kelvin}$.
\end{talist}

\begin{solution}
\begin{talist}
\ta $\Delta V_{BE} = \SI{25.865}{\milli\volt}\cdot\ln 8 = \SI{53.8}{\milli\volt}$,
	$I = \SI{53.8}{\milli\volt}/\SI{1}{\kilo\ohm} = \SI{53.8}{\micro\ampere}$,
	$V_E = 2\cdot\SI{53.8}{\micro\ampere}\cdot\SI{5.6}{\kilo\ohm} = \SI{0.602}{\volt}$.
\ta $V_{BG} = \SI{0.520}{\volt}+\SI{0.602}{\volt} = \SI{1.122}{\volt}$.
\ta $\mathrm{d}V_E/\mathrm{d}T = 2\frac{R_1}{R_2}\frac{k}{q}\ln 8 = 11.2\cdot\SI{86.2}{\micro\volt/\kelvin}\cdot 2.08 = +\SI{2.0}{\milli\volt/\kelvin}$:
	it cancels the TC of $V_{BE}$.
\item[] \falstad{brokaw_cell}: \texttt{VBG} $=\SI{1.122}{\volt}$, \texttt{VE} $=\SI{0.602}{\volt}$;
	\texttt{C1} and \texttt{C2} are equal (\SI{4.47}{\volt}), the voltage across $R_2$ is \SI{53.8}{\milli\volt}.
\end{talist}
\end{solution}

%%%%%%%%%%%%%%%%%%%%%%%%%%%%%%%%%%%%%%%%%%%%%%%%%%%%%%%%%%%%%%%%%%%%%%%%%%%
\section{Shunt and Series References}
%%%%%%%%%%%%%%%%%%%%%%%%%%%%%%%%%%%%%%%%%%%%%%%%%%%%%%%%%%%%%%%%%%%%%%%%%%%
\nsl{Integrated references come in two kinds. A \emph{shunt} reference
behaves like a precise Zener diode: it is connected in parallel to the load
and draws exactly the current needed to keep its voltage constant; it is
fed through a series resistor. A \emph{series} reference behaves like a
precise voltage regulator: it has an input, a ground and an output pin, and
it only draws a small quiescent current itself. The most widely used shunt
reference is the TL431, an adjustable bandgap reference with an internal
\SI{2.495}{\volt} reference, an error amplifier and a shunt transistor.\newpage}

\begin{description}
\item[{\rot\bf The adjustable shunt reference TL431}]~\\[-\lblskip]
\begin{center}
\begin{circuitikz}[scale=0.8, transform shape]
  \draw (0,0) to[V, invert, l=$V_S$] (0,4) to[R, l=$R_V$] (3,4) -- (5.5,4);
  \draw (3,4) to[R, l_=$R_1$] (3,2) to[R, l_=$R_2$] (3,0);
  \draw (5,4) to[zD, invert, n=tl] (5,0);
  \draw (3,2) node[circ]{} -- (4.4,2) -- (4.85,2.15);
  \node[right,font=\scriptsize] at (5.3,2.4) {TL431};
  \node[font=\scriptsize] at (3.8,2.25) {REF};
  \node[right,font=\scriptsize] at (5.1,3.6) {K};
  \node[right,font=\scriptsize] at (5.1,0.4) {A};
  \draw (5.5,4) to[short,-o] (7,4) node[right]{$\Vout$};
  \draw (0,0) -- (5,0);
  \draw (0,0) node[ground]{};
  \draw (3,4) node[circ]{} (5,4) node[circ]{};
\end{circuitikz}
\end{center}
\bi
	\ite the shunt element conducts until $V_\mathrm{REF} = \SI{2.5}{\volt}$
	\keyeq{TL431}{$\displaystyle \Vout = V_\mathrm{REF}\left(1+\frac{R_1}{R_2}\right) + I_\mathrm{REF}\,R_1$, \quad $V_\mathrm{REF} = \SI{2.495}{\volt}$, $I_\mathrm{REF}\approx\SI{2}{\micro\ampere}$}
	\ite cathode current between \SI{1}{\milli\ampere} and \SI{100}{\milli\ampere}
\ei
\os{\newpage}

\item[{\rot\bf Inside the TL431 (behavioural model)}]~\\[-\lblskip]
\begin{center}
\begin{circuitikz}[scale=0.85, transform shape]
  \coordinate (gnd) at (0,0);
  \node[op amp, noinv input up] (A) at (2.9,2) {};
  \draw (A.+) -- ++(-0.9,0) to[short,-o] (-0.4,0 |- A.+) node[left]{REF};
  \draw (A.-) -- ++(-0.5,0) coordinate(m);
  \draw (m |- gnd) node[circ]{} to[V, invert, l=\SI{2.5}{\volt}] (m);
  \node[npn] (Q) at (6.4,2) {};
  \draw (A.out) to[R, l=\SI{10}{\kilo\ohm}] (Q.B);
  \draw (-0.4,4) node[ocirc]{} node[left]{K} -- (Q.C |- 0,4) -- (Q.C);
  \draw (-0.4,0) node[ocirc]{} node[left]{A} -- (Q.E |- gnd) -- (Q.E);
  \draw[dashed,rwugray] (0.0,-0.3) rectangle (7.2,4.3);
\end{circuitikz}
\end{center}
\bi
	\ite error amplifier compares REF with the internal \SI{2.5}{\volt} bandgap
	\ite $V_\mathrm{REF}$ too high $\Rightarrow$ transistor conducts more $\Rightarrow$ $\Vout$ falls
	\ite $R_1 = \SI{12}{\kilo\ohm}$, $R_2 = \SI{20}{\kilo\ohm}$: $\Vout = 2.5\cdot 1.6 = \SI{4.00}{\volt}$
\ei
\falstad{shunt_ref}
\nsl{\par The Falstad model sets the TL431 to \SI{4.00}{\volt}, fed from
\SI{12}{\volt} through \SI{1}{\kilo\ohm}. The cathode current is
$(\SI{12}{\volt}-\SI{4}{\volt})/\SI{1}{\kilo\ohm} - \SI{4}{\volt}/\SI{32}{\kilo\ohm}
= \SI{7.875}{\milli\ampere}$. Move the supply slider down to \SI{6}{\volt}:
the output stays at \SI{4.00}{\volt}, only the cathode current falls. The
internal reference is modelled as an ideal \SI{2.5}{\volt} source; the real
chip uses a bandgap cell.}
\os{\newpage}

\item[{\rot\bf Shunt or series reference?}]~\\[-\lblskip]
\bi
	\ite \textbf{shunt} (TL431, LM4040, LM385)
	\itee two terminals, works at any supply above $\Vout$, can be used as a negative reference
	\itee the series resistor wastes current: $R_V$ must deliver $I_\mathrm{min} + I_{L,\max}$ at the lowest supply
	\ite \textbf{series} (REF50xx, ADR4xx, MAX6xxx)
	\itee three terminals, quiescent current only \SIrange{0.1}{1}{\milli\ampere}
	\itee better for varying loads and battery operation, needs $V_\mathrm{in} > \Vout + V_\mathrm{dropout}$
	\ite both: decouple according to the data sheet (some are unstable with the wrong capacitor)
\ei
\os{\newpage}
\end{description}

\exercise{TL431 set to 4.00 V}
A TL431 ($V_\mathrm{REF} = \SI{2.495}{\volt}$, $I_\mathrm{REF} = \SI{2}{\micro\ampere}$,
$I_{K,\min} = \SI{1}{\milli\ampere}$) shall provide \SI{4.00}{\volt} from a
supply of \SIrange{6}{15}{\volt}.
\begin{talist}
\ta Choose $R_1$ and $R_2$ from the E24 series.
\ta How large is the error caused by $I_\mathrm{REF}$ and by $V_\mathrm{REF} = \SI{2.495}{\volt}$ instead of \SI{2.5}{\volt}?
\ta Choose $R_V$ such that a load of \SI{1}{\milli\ampere} can be supplied at \SI{6}{\volt}.
\end{talist}

\begin{solution}
\begin{talist}
\ta $1+R_1/R_2 = 1.6 \Rightarrow R_1/R_2 = 0.6$: $R_1 = \SI{12}{\kilo\ohm}$, $R_2 = \SI{20}{\kilo\ohm}$ (both E24).
\ta $I_\mathrm{REF}R_1 = \SI{2}{\micro\ampere}\cdot\SI{12}{\kilo\ohm} = \SI{24}{\milli\volt}$ ($+\SI{0.6}{\percent}$);
	$\SI{2.495}{\volt}\cdot 1.6 = \SI{3.992}{\volt}$ ($-\SI{0.2}{\percent}$). Together $\Vout = \SI{4.016}{\volt}$:
	for an accurate \SI{4.00}{\volt} choose smaller divider resistors or trim.
\ta At \SI{6}{\volt} $R_V$ must carry $I_{K,\min} + I_L + I_\mathrm{div} = 1 + 1 + 0.125 = \SI{2.125}{\milli\ampere}$:
	$R_V \le (\SI{6}{\volt}-\SI{4}{\volt})/\SI{2.125}{\milli\ampere} = \SI{941}{\ohm}$, choose \SI{910}{\ohm}.
	At \SI{15}{\volt} the TL431 then carries $(15-4)/0.91 - 0.125 - 1 = \SI{11}{\milli\ampere}$ (\SI{44}{\milli\watt}): fine.
\item[] \falstad{shunt_ref} (ideal model with \SI{2.5}{\volt}, $R_V = \SI{1}{\kilo\ohm}$, no load):
	\texttt{VK} $=\SI{4.00}{\volt}$ at \SI{12}{\volt} and at \SI{6}{\volt}, \texttt{REF} $= \SI{2.50}{\volt}$.
\end{talist}
\end{solution}

%%%%%%%%%%%%%%%%%%%%%%%%%%%%%%%%%%%%%%%%%%%%%%%%%%%%%%%%%%%%%%%%%%%%%%%%%%%
\section{Buffering and Scaling a Reference}
%%%%%%%%%%%%%%%%%%%%%%%%%%%%%%%%%%%%%%%%%%%%%%%%%%%%%%%%%%%%%%%%%%%%%%%%%%%
\nsl{References are available in a few standard voltages: \SI{1.25}{},
\SI{2.048}{}, \SI{2.5}{}, \SI{3.0}{}, \SI{4.096}{} and \SI{5.0}{\volt}.
If the circuit needs a different value, or more current than the reference
can deliver, an op-amp buffers and scales it. The capstone needs exactly
\SI{4.00}{\volt} for a reference ladder of eight \SI{1}{\kilo\ohm}
resistors, so that one LSB of the 3-bit flash converter is
\SI{0.5}{\volt}. The ladder draws $\SI{4}{\volt}/\SI{8}{\kilo\ohm} =
\SI{0.5}{\milli\ampere}$: far too much for a voltage divider with
high-value resistors, and a noticeable load for many references. The
op-amp circuits are treated in detail in chapter~5; here we only use the
results.\newpage}

\begin{description}
\item[{\rot\bf A divider is no reference}]~\\[-\lblskip]
\begin{center}
\begin{circuitikz}[scale=0.8, transform shape]
  \draw (0,0) to[V, invert, l=\SI{5.00}{\volt}] (0,3) -- (1.5,3)
        to[R, l=\SI{10}{\kilo\ohm}] (1.5,1.5) to[R, l=\SI{40}{\kilo\ohm}] (1.5,0) -- (0,0);
  \draw (1.5,1.5) node[circ]{} -- (3.9,1.5) to[R, l=\SI{8}{\kilo\ohm}] (3.9,0) -- (1.5,0);
  \draw (0,0) node[ground]{};
  \node[font=\small] at (2.2,-0.8) {$V = \SI{2.00}{\volt}$};
  \draw (6,0) to[V, invert, l=\SI{5.00}{\volt}] (6,3) -- (7.5,3)
        to[R, l=\SI{10}{\kilo\ohm}] (7.5,1.5) to[R, l=\SI{40}{\kilo\ohm}] (7.5,0) -- (6,0);
  \draw (6,0) node[ground]{};
  \node[op amp, noinv input up, anchor=+] (A) at (9.4,1.5) {};
  \draw (7.5,1.5) node[circ]{} -- (A.+);
  \draw (A.-) -- ++(0,-0.6) -| (A.out);
  \draw (A.out) -- ++(0.6,0) coordinate(o) to[R, l=\SI{8}{\kilo\ohm}] (o |- 0,-0.2) node[ground]{};
  \node[font=\small] at (10.2,-0.8) {$V = \SI{4.00}{\volt}$};
\end{circuitikz}
\end{center}
\bi
	\ite divider $\SI{10}{\kilo\ohm}/\SI{40}{\kilo\ohm}$: \SI{4.00}{\volt} open, but only \SI{2.00}{\volt} with the \SI{8}{\kilo\ohm} ladder
	\ite source resistance $R_1\parallelto R_2 = \SI{8}{\kilo\ohm}$ = load resistance
	\ite follower: input current $\approx 0$, output resistance $\approx 0$: \SI{4.00}{\volt}
\ei
\falstad{ref_divider_loaded}
\os{\newpage}

\item[{\rot\bf Scaling with a non-inverting amplifier}]~\\[-\lblskip]
\begin{center}
\begin{circuitikz}[scale=0.85, transform shape]
  \draw (-2.8,0) to[V, invert, l=\SI{2.5}{\volt}] (-2.8,2.6) -- (0.3,2.6);
  \draw (-2.8,0) node[ground]{};
  \node[op amp, noinv input up, anchor=+] (A) at (0.3,2.6) {};
  \draw (A.-) -- ++(-0.3,0) -- ++(0,-0.9) coordinate(m);
  \draw (m) to[R, l_=$R_G{=}\SI{20}{\kilo\ohm}$] (m |- 0,-0.2) node[ground]{};
  \draw (m) node[circ]{} to[R, l_=$R_F{=}\SI{12}{\kilo\ohm}$] (3.4,0 |- m) -- (3.4,0 |- A.out);
  \draw (A.out) -- (3.4,0 |- A.out) node[circ]{} -- (5,0 |- A.out) node[circ]{} node[above]{\SI{4.00}{\volt}}
        to[R] ++(0,-1) to[R] ++(0,-1) -- ++(0,-0.2) node[below]{$\vdots$};
  \node[right,font=\small] at (5.3,1.2) {ladder $8\times\SI{1}{\kilo\ohm}$};
\end{circuitikz}
\end{center}
\bi
	\ite $G = 1+R_F/R_G = 1 + 12/20 = 1.6$: \quad $\SI{2.5}{\volt}\cdot 1.6 = \SI{4.00}{\volt}$
	\ite the op-amp delivers ladder current \SI{0.5}{\milli\ampere} plus \SI{0.125}{\milli\ampere} for $R_F+R_G$
	\ite gain error from resistor tolerance: $\frac{\Delta G}{G} = \frac{G-1}{G}\left(\frac{\Delta R_F}{R_F}-\frac{\Delta R_G}{R_G}\right)$
\ei
\falstad{ref_buffer_ladder}
\nsl{\par The resistor tolerances enter the result only with the factor
$(G-1)/G = 0.375$: two \SI{0.1}{\percent} resistors give at most
$0.375\cdot\SI{0.2}{\percent} = \SI{0.075}{\percent}$. The offset voltage
of the op-amp is multiplied by $G$: \SI{1}{\milli\volt} offset gives
\SI{1.6}{\milli\volt} at the output (\SI{0.04}{\percent}). Alternatively a
\SI{4.096}{\volt} reference with a follower can be used; the LSB is then
\SI{0.512}{\volt} instead of \SI{0.5}{\volt}.}
\os{\newpage}

\item[{\rot\bf Design card}]~\\[-\lblskip]
\designcard{Reference for an ADC ladder}{
\bi
	\ite full scale $V_\mathrm{FS}$ and LSB: $\LSB = V_\mathrm{FS}/2^N$ \quad (\SI{4.00}{\volt}, 3 bit: \SI{0.5}{\volt})
	\ite ladder current $I = V_\mathrm{FS}/(2^N R)$ \quad ($8\times\SI{1}{\kilo\ohm}$: \SI{0.5}{\milli\ampere})
	\ite reference voltage $V_R$ from a standard value, gain $G = V_\mathrm{FS}/V_R$ \quad (2.5 V, $G = 1.6$)
	\ite non-inverting amplifier: $R_F/R_G = G-1$, $R_F + R_G \gg$ ladder ($\SI{12}{\kilo\ohm}+\SI{20}{\kilo\ohm}$)
	\ite check: op-amp output current, $(G-1)/G\cdot$ resistor tolerance, $G\cdot V_\mathrm{offset}$
	\ite decouple the ladder top with \SIrange{100}{1000}{\nano\farad} close to the comparators
\ei}
\os{\newpage}
\end{description}

\exercise{Reference for the 3-bit flash converter}
The flash converter of the capstone needs \SI{4.00}{\volt} at the top of a
ladder of $8\times\SI{1}{\kilo\ohm}$. A \SI{2.5}{\volt} reference is available.
\begin{talist}
\ta Calculate the ladder current and the tap voltages.
\ta Design a non-inverting amplifier with E24 resistors.
\ta Why can the \SI{4.00}{\volt} not be taken from a divider \SI{10}{\kilo\ohm}/\SI{40}{\kilo\ohm} at a \SI{5}{\volt} reference?
\ta How large is the worst-case error of $\Vref$ with \SI{1}{\percent} resistors?
\end{talist}

\begin{solution}
\begin{talist}
\ta $I = \SI{4}{\volt}/\SI{8}{\kilo\ohm} = \SI{0.5}{\milli\ampere}$; taps every \SI{0.5}{\volt}: $T_1 = \SI{0.5}{\volt}$, \ldots, $T_4 = \SI{2.0}{\volt}$, \ldots, $T_7 = \SI{3.5}{\volt}$.
\ta $G = 4/2.5 = 1.6$, $R_F/R_G = 0.6$: $R_F = \SI{12}{\kilo\ohm}$, $R_G = \SI{20}{\kilo\ohm}$.
\ta The divider has a source resistance of $\SI{10}{\kilo\ohm}\parallelto\SI{40}{\kilo\ohm} = \SI{8}{\kilo\ohm}$, equal to the
	ladder. The loaded voltage is $5\cdot\frac{40\parallelto 8}{10 + 40\parallelto 8} = 5\cdot\frac{6.67}{16.67} = \SI{2.00}{\volt}$.
\ta $\Delta G/G = 0.375\cdot 2\cdot\SI{1}{\percent} = \SI{0.75}{\percent}$, i.e.\ $\pm\SI{30}{\milli\volt}$;
	for the 3-bit converter (LSB \SI{0.5}{\volt}) irrelevant, for a 12-bit converter 30 LSB.
\item[] \falstad{ref_buffer_ladder}: \texttt{VREF4} $=\SI{4.00}{\volt}$, \texttt{T4} $=\SI{2.00}{\volt}$,
	\texttt{T1} $=\SI{0.50}{\volt}$, ladder current \SI{0.5}{\milli\ampere};
	\falstad{ref_divider_loaded}: \texttt{Vopen} $=\SI{4.00}{\volt}$, \texttt{Vload} $=\SI{2.00}{\volt}$,
	\texttt{Vbuf} $=\SI{4.00}{\volt}$.
\end{talist}
\end{solution}

%%%%%%%%%%%%%%%%%%%%%%%%%%%%%%%%%%%%%%%%%%%%%%%%%%%%%%%%%%%%%%%%%%%%%%%%%%%
\section{Ratiometric Measurement}
%%%%%%%%%%%%%%%%%%%%%%%%%%%%%%%%%%%%%%%%%%%%%%%%%%%%%%%%%%%%%%%%%%%%%%%%%%%
\nsl{Resistive sensors do not produce a voltage themselves; they need an
excitation, and their output is a fraction of it. If the same reference
excites the sensor and sets the full scale of the ADC, the reference
appears in the numerator and in the denominator of the result and cancels.
Such a measurement is called \emph{ratiometric}. It is the most effective
way to get an accurate result from an inexpensive reference: the reference
only needs to be quiet and stable during one conversion, not accurate. The
capstone uses this principle: the \SI{4.00}{\volt} reference feeds both the
sensor bridge and the ladder of the flash converter \cite{kester2005}.\newpage}

\begin{description}
\item[{\rot\bf The reference cancels}]~\\[-\lblskip]
\begin{center}
\begin{circuitikz}[scale=0.8, transform shape]
  \draw (0,4) node[above]{$\Vref$} -- (3.5,4);
  \draw (0,4) to[R, l_=$R_1$] (0,2) to[vR, l_=$R_S$] (0,0) node[ground]{};
  \draw (0,2) node[circ]{} -- (1.8,2) node[right]{$\Vin = \Vref\,\frac{R_S}{R_1+R_S}$};
  \draw (3.5,4) -- (6.5,4) to[R] (6.5,3) to[R] (6.5,2) to[R] (6.5,1) to[R] (6.5,0) node[ground]{};
  \node[right,font=\small,align=left] at (6.7,2) {ladder:\\thresholds\\$\frac{k}{4}\Vref$};
\end{circuitikz}
\end{center}
\keyeq{Ratiometric}{$\displaystyle \mathrm{code} = 2^N\frac{\Vin}{\Vref} = 2^N\frac{R_S}{R_1+R_S}$ \qquad independent of $\Vref$}
\os{\newpage}

\item[{\rot\bf Ratiometric versus absolute}]~\\[-\lblskip]
\bi
	\ite Falstad: two 2-bit flash converters (3 comparators), sensor divider \SI{1}{\kilo\ohm}/\SI{2.7}{\kilo\ohm}: $\Vin = 0.730\,V_\mathrm{exc}$
	\ite left: sensor and ladder from $\Vref$: code 2 for $\Vref = $ \SI{3.6}{}, \SI{4.0}{}, \SI{4.4}{\volt}
	\ite right: sensor from a fixed \SI{4.00}{\volt}, ladder from $\Vref$
	\itee $\Vref = \SI{4.0}{\volt}$: $\Vin = \SI{2.92}{\volt} < \SI{3.0}{\volt}$: code 2
	\itee $\Vref = \SI{3.6}{\volt}$: threshold $\frac34\Vref = \SI{2.7}{\volt}$: code 3 (wrong)
	\ite ratiometric works only if both see the \emph{same} voltage at the \emph{same} time
	\itee a common connection with voltage drops, or noise on one path only, breaks it
\ei
\falstad{ratiometric}
\os{\newpage}
\end{description}

%%%%%%%%%%%%%%%%%%%%%%%%%%%%%%%%%%%%%%%%%%%%%%%%%%%%%%%%%%%%%%%%%%%%%%%%%%%
\section{Error Budget}
%%%%%%%%%%%%%%%%%%%%%%%%%%%%%%%%%%%%%%%%%%%%%%%%%%%%%%%%%%%%%%%%%%%%%%%%%%%
\nsl{An error budget lists all contributions to the error of a quantity
and adds them up. Independent random errors are usually added as the root
of the sum of squares (RSS), which is the realistic expectation;
the simple sum is the worst case, which a guaranteed specification must
meet. We look at the \SI{4.00}{\volt} reference of the capstone, built from
a \SI{2.5}{\volt} series reference and the non-inverting amplifier, over an
ambient temperature of \SIrange{0}{50}{\celsius}.\newpage}

\begin{description}
\item[{\rot\bf Error budget of the 4.00\,V reference}]~\\[-\lblskip]
\begin{center}
\begin{tabular}{lll}
\toprule
contribution & data & error \\
\midrule
initial accuracy & $\pm\SI{0.2}{\percent}$ & \SI{0.20}{\percent} \\
temperature & \SI{25}{ppm/K}, $\pm\SI{25}{\kelvin}$ & \SI{0.06}{\percent} \\
gain resistors & 2 $\times$ \SI{0.1}{\percent}, factor 0.375 & \SI{0.08}{\percent} \\
op-amp offset & \SI{1}{\milli\volt} $\times$ 1.6 & \SI{0.04}{\percent} \\
long-term drift & \SI{50}{ppm} / \SI{1000}{\hour} & \SI{0.01}{\percent} \\
\midrule
worst case (sum) & & \SI{0.39}{\percent} \\
RSS & & \SI{0.23}{\percent} \\
\bottomrule
\end{tabular}
\end{center}
\bi
	\ite compare: 3 bit, $\LSB = \SI{12.5}{\percent}$ of full scale: no problem at all
	\ite 12 bit, $\LSB = \SI{0.024}{\percent}$: 16 LSB worst case $\Rightarrow$ calibrate or measure ratiometrically
\ei
\nsl{The initial accuracy dominates. It can be removed by a one-time
calibration (trimming, or a correction factor in the software). What
remains after calibration is temperature drift and ageing, and these are
the parameters that distinguish an expensive reference from a cheap one.
In a ratiometric design all reference terms drop out; only the ratio
accuracy of the resistors remains.}
\os{\newpage}
\end{description}

\exercise{Ratiometric measurement and error budget}
A sensor divider $R_1 = \SI{1}{\kilo\ohm}$, $R_S = \SI{2.7}{\kilo\ohm}$ is read
by a 2-bit flash converter with a ladder of $4\times\SI{1}{\kilo\ohm}$.
\begin{talist}
\ta The divider and the ladder are fed from the same $\Vref$. Which code results for $\Vref = \SI{4.0}{\volt}$ and \SI{3.6}{\volt}?
\ta The divider is fed from a fixed \SI{4.00}{\volt}, the ladder from $\Vref$. Repeat a).
\ta A 12-bit converter uses a reference with $\pm\SI{0.2}{\percent}$ initial accuracy and \SI{25}{ppm/K}.
	How many LSB can the result be off at \SI{50}{\celsius} after calibration at \SI{25}{\celsius}?
\end{talist}

\begin{solution}
\begin{talist}
\ta $\Vin/\Vref = 2.7/3.7 = 0.730$; the thresholds are at $0.25$, $0.5$ and $0.75\,\Vref$.
	$0.5 < 0.730 < 0.75$: code 2 for every $\Vref$.
\ta $\Vin = \SI{2.92}{\volt}$. $\Vref = \SI{4.0}{\volt}$: thresholds 1, 2, \SI{3}{\volt}: code 2.
	$\Vref = \SI{3.6}{\volt}$: thresholds 0.9, 1.8, \SI{2.7}{\volt}: code 3.
	A \SI{10}{\percent} reference error changes the result by one code.
\ta After calibration only the drift remains: $\SI{25}{ppm/K}\cdot\SI{25}{\kelvin} = \SI{625}{ppm}$
	$= 0.000625\cdot 4096 = 2.6$ LSB. Without calibration add $0.002\cdot 4096 = 8.2$ LSB.
\item[] \falstad{ratiometric}: with the \texttt{VREF} slider at \SI{4.0}{\volt} both converters show
	D = E = code 2 (\texttt{VIN1} $=\SI{2.92}{\volt}$); at \SI{3.6}{\volt} the right converter switches \texttt{E3} on (code 3), the left stays at code 2.
\end{talist}
\end{solution}

%%% Local Variables:
%%% mode: latex
%%% TeX-master: "docu"
%%% End:
